\documentclass[a4paper,11pt]{article}

\usepackage[utf8]{inputenc}
\usepackage[T1]{fontenc}
\usepackage[english]{babel}
\usepackage{amsmath,amssymb,amsfonts}
\providecommand{\ket}[1]{\mathinner{|{#1}\rangle}}
\providecommand{\bra}[1]{\mathinner{\langle{#1}|}}
\providecommand{\braket}[2]{\mathinner{\langle{#1}|\,{#2}\rangle}}
\usepackage{geometry}
\geometry{margin=2.5cm}

\title{Quantum Interference in Quantum Computing:\\ A Mathematical Introduction}
\author{}
\date{}

\begin{document}

\maketitle

\section*{Introduction}

\emph{Quantum interference} is one of the most fundamental phenomena that distinguishes quantum
computation from classical computation. It refers to the way probability amplitudes---complex
numbers associated with quantum states---combine according to the rules of wave mechanics:
they can \emph{constructively interfere} (reinforcing each other) or \emph{destructively interfere}
(canceling each other out).

In quantum algorithms, interference is carefully engineered through sequences of quantum gates
to amplify the probability of correct answers while suppressing incorrect ones. This document
presents a mathematical introduction to quantum interference using simple one- and two-qubit
examples, and then connects these ideas to Grover's search algorithm.

\section*{Classical vs.\ Quantum Probability}

In classical probability theory, probabilities are non-negative real numbers that add linearly.
If an event can occur via two mutually exclusive paths with probabilities $p_1$ and $p_2$, the
total probability is
\begin{equation}
  P_{\text{total}} = p_1 + p_2.
\end{equation}

In quantum mechanics, each computational path is associated with a \emph{complex probability
amplitude} $\alpha \in \mathbb{C}$. The probability of an outcome is given by the squared modulus
of the \emph{sum} of amplitudes:
\begin{equation}
  P_{\text{quantum}} = |\alpha_1 + \alpha_2|^2.
\end{equation}

Because amplitudes are complex, they have both magnitude and phase. When two amplitudes are
added, their phases determine whether they reinforce (constructive interference) or cancel
(destructive interference). This is fundamentally different from classical probability addition.

\section*{Mathematical Preliminaries}

\subsection*{Qubit States and Superposition}

A single qubit lives in a two-dimensional Hilbert space spanned by the computational basis states
$\ket{0}$ and $\ket{1}$. A general qubit state is a \emph{superposition}:
\begin{equation}
  \ket{\psi} = \alpha\ket{0} + \beta\ket{1},
\end{equation}
where $\alpha, \beta \in \mathbb{C}$ satisfy the normalization condition
\begin{equation}
  |\alpha|^2 + |\beta|^2 = 1.
\end{equation}

Upon measurement in the computational basis, the qubit collapses to $\ket{0}$ with probability
$|\alpha|^2$ and to $\ket{1}$ with probability $|\beta|^2$.

\subsection*{The Hadamard Gate}

The Hadamard gate $H$ is defined by
\begin{equation}
  H\ket{0} = \frac{1}{\sqrt{2}}\bigl(\ket{0} + \ket{1}\bigr) \equiv \ket{+},
  \qquad
  H\ket{1} = \frac{1}{\sqrt{2}}\bigl(\ket{0} - \ket{1}\bigr) \equiv \ket{-}.
\end{equation}
In matrix form,
\begin{equation}
  H = \frac{1}{\sqrt{2}}
  \begin{pmatrix}
    1 & 1 \\
    1 & -1
  \end{pmatrix},
  \qquad H^2 = I.
\end{equation}

\section*{Example 1: Interference in the $H$--$H$ Circuit}

Consider the circuit that applies two Hadamard gates in succession to a qubit initialized in
$\ket{0}$.

\subsection*{Step 1: First Hadamard creates superposition}

From $\ket{0}$, the first Hadamard gate produces
\begin{equation}
  \ket{\psi_1} = H\ket{0} = \frac{1}{\sqrt{2}}\bigl(\ket{0} + \ket{1}\bigr) = \ket{+}.
\end{equation}

At this stage, the qubit is in an equal superposition of both computational basis states. If we
measured now, we would obtain $\ket{0}$ or $\ket{1}$ each with probability $\frac{1}{2}$.

\subsection*{Step 2: Second Hadamard causes interference}

We now apply a second Hadamard:
\begin{align}
  \ket{\psi_2} &= H\ket{+} = H\left(\frac{1}{\sqrt{2}}\bigl(\ket{0} + \ket{1}\bigr)\right) \\
  &= \frac{1}{\sqrt{2}}\bigl(H\ket{0} + H\ket{1}\bigr) \\
  &= \frac{1}{\sqrt{2}}\left(\frac{1}{\sqrt{2}}\bigl(\ket{0} + \ket{1}\bigr)
     + \frac{1}{\sqrt{2}}\bigl(\ket{0} - \ket{1}\bigr)\right) \\
  &= \frac{1}{2}\Bigl[\bigl(\ket{0} + \ket{1}\bigr) + \bigl(\ket{0} - \ket{1}\bigr)\Bigr] \\
  &= \frac{1}{2}\bigl(2\ket{0}\bigr) = \ket{0}.
\end{align}

Thus the final state is $\ket{0}$ with certainty. The amplitudes for $\ket{1}$ have interfered
destructively.

\subsection*{Path Interpretation}

After the first Hadamard, the qubit is in
\[
  \frac{1}{\sqrt{2}}\ket{0} + \frac{1}{\sqrt{2}}\ket{1}.
\]
When the second Hadamard is applied, the contribution to $\ket{1}$ from the ``path via $\ket{0}$''
is $+\frac{1}{2}$, while the contribution from the ``path via $\ket{1}$'' is $-\frac{1}{2}$. These
cancel, giving zero total amplitude for $\ket{1}$. This is destructive interference.

\section*{Example 2: Phase Flip and Interference Pattern}

Consider the circuit $\ket{0} \xrightarrow{H} \xrightarrow{Z} \xrightarrow{H}$, where the
Pauli-$Z$ gate is
\begin{equation}
  Z =
  \begin{pmatrix}
    1 & 0 \\
    0 & -1
  \end{pmatrix},
  \qquad Z\ket{0} = \ket{0}, \quad Z\ket{1} = -\ket{1}.
\end{equation}

\subsection*{Step-by-step evolution}

\textbf{First Hadamard:}
\begin{equation}
  \ket{\psi_1} = H\ket{0} = \frac{1}{\sqrt{2}}\bigl(\ket{0} + \ket{1}\bigr).
\end{equation}

\textbf{Phase flip via $Z$:}
\begin{align}
  \ket{\psi_2} &= Z\ket{\psi_1} = \frac{1}{\sqrt{2}}\bigl(Z\ket{0} + Z\ket{1}\bigr) \\
  &= \frac{1}{\sqrt{2}}\bigl(\ket{0} - \ket{1}\bigr) = \ket{-}.
\end{align}

\textbf{Second Hadamard:}
\begin{align}
  \ket{\psi_3} &= H\ket{-} = H\left(\frac{\ket{0} - \ket{1}}{\sqrt{2}}\right) \\
  &= \frac{1}{\sqrt{2}}\bigl(H\ket{0} - H\ket{1}\bigr) \\
  &= \frac{1}{\sqrt{2}}\left(\frac{\ket{0} + \ket{1}}{\sqrt{2}}
     - \frac{\ket{0} - \ket{1}}{\sqrt{2}}\right) \\
  &= \frac{1}{2}\Bigl[(\ket{0} + \ket{1}) - (\ket{0} - \ket{1})\Bigr] \\
  &= \frac{1}{2}(2\ket{1}) = \ket{1}.
\end{align}

The $Z$ gate changed the relative phase, thus flipping the interference pattern: now $\ket{1}$
experiences constructive interference and $\ket{0}$ destructive interference.

\section*{Example 3: Two-Qubit Interference (Deutsch--Jozsa with $n = 2$)}

We now consider a simple two-qubit example that generalises the one-qubit interference picture
and is closely related to the Deutsch--Jozsa algorithm.

\subsection*{Setup}

Consider two input qubits and one ancilla qubit, all initialised as
\begin{equation}
  \ket{\psi_0} = \ket{00}\ket{1}.
\end{equation}
We will:
\begin{enumerate}
  \item Apply Hadamards to all three qubits.
  \item Apply an oracle $U_f$ that encodes a Boolean function $f : \{0,1\}^2 \to \{0,1\}$.
  \item Apply Hadamards again to the two input qubits.
  \item Measure the input register.
\end{enumerate}

The oracle is defined by
\begin{equation}
  U_f\ket{x}\ket{y} = \ket{x}\ket{y \oplus f(x)},
\end{equation}
where $x \in \{0,1\}^2$ and $\oplus$ is XOR.

\subsection*{Step 1: Create uniform superposition}

Applying $H^{\otimes 3}$ to $\ket{00}\ket{1}$ yields
\begin{align}
  \ket{\psi_1} &= (H^{\otimes 2}\ket{00}) \otimes (H\ket{1}) \\
  &= \left(\frac{1}{2}\sum_{x \in \{0,1\}^2}\ket{x}\right)
     \otimes \left(\frac{\ket{0} - \ket{1}}{\sqrt{2}}\right) \\
  &= \frac{1}{2\sqrt{2}}\sum_{x \in \{0,1\}^2}\ket{x}\bigl(\ket{0} - \ket{1}\bigr).
\end{align}

Thus each input basis state $\ket{x}$ appears with equal amplitude, and the ancilla is in the
$\ket{-}$ state.

\subsection*{Step 2: Apply the oracle $U_f$}

We act with $U_f$:
\begin{align}
  \ket{\psi_2} &= \frac{1}{2\sqrt{2}}\sum_{x \in \{0,1\}^2}
     U_f\Bigl(\ket{x}\bigl(\ket{0} - \ket{1}\bigr)\Bigr) \\
  &= \frac{1}{2\sqrt{2}}\sum_{x}
     \left(\ket{x}\ket{0 \oplus f(x)} - \ket{x}\ket{1 \oplus f(x)}\right).
\end{align}

We use the fact that
\[
  \ket{0 \oplus f(x)} - \ket{1 \oplus f(x)} = (-1)^{f(x)}\bigl(\ket{0} - \ket{1}\bigr),
\]
so
\begin{align}
  \ket{\psi_2} &= \frac{1}{2\sqrt{2}}\sum_{x}(-1)^{f(x)}\ket{x}\bigl(\ket{0} - \ket{1}\bigr) \\
  &= \left(\frac{1}{2}\sum_{x}(-1)^{f(x)}\ket{x}\right) \otimes \ket{-}.
\end{align}

The ancilla returns to $\ket{-}$ and all information about $f$ is now encoded in the phases
$(-1)^{f(x)}$ of the input superposition. This is a phase kickback effect.

\subsection*{Step 3: Final Hadamards and interference}

We now apply $H^{\otimes 2}$ to the input qubits:
\begin{equation}
  \ket{\psi_3} = \left(H^{\otimes 2}\frac{1}{2}\sum_{x}(-1)^{f(x)}\ket{x}\right) \otimes \ket{-}.
\end{equation}

Using
\[
  H^{\otimes 2}\ket{x} = \frac{1}{2}\sum_{z \in \{0,1\}^2}(-1)^{x \cdot z}\ket{z},
\]
where $x \cdot z$ is the bitwise inner product modulo 2, we obtain
\begin{equation}
  \ket{\psi_3} = \frac{1}{4}\sum_{x,z}(-1)^{f(x)}(-1)^{x \cdot z}\ket{z} \otimes \ket{-}.
\end{equation}

The amplitude of each output basis state $\ket{z}$ is
\[
  \alpha_z = \frac{1}{4}\sum_{x}(-1)^{f(x) + x \cdot z}.
\]

\paragraph{Constant $f$.}
If $f$ is constant (always $0$ or always $1$), the phase $(-1)^{f(x)}$ is the same for all $x$. In
that case, for $z = 00$ every term in the sum has the same sign and adds constructively, while for
$z \neq 00$ the terms cancel pairwise. Thus the input register ends up in $\ket{00}$ with
certainty.

\paragraph{Balanced $f$.}
If $f$ is balanced (half of the inputs map to $0$ and half to $1$), the sum for $z = 00$ splits into
two equal halves with opposite signs and cancels, so the amplitude of $\ket{00}$ is zero. The final
state has no support on $\ket{00}$; interference has removed it completely.

This two-qubit example shows how phase information introduced by the oracle is converted, via
interference, into observable differences in measurement outcomes.

\section*{Interference in Grover's Algorithm}

We now briefly describe how interference appears in Grover's search algorithm through the
\emph{diffusion operator}.

\subsection*{Uniform superposition and marking}

Consider an $n$-qubit register with computational basis $\{\ket{x}\}_{x \in \{0,1\}^n}$. Start from
\begin{equation}
  \ket{\psi_0} = \ket{0}^{\otimes n},
\end{equation}
then apply $H^{\otimes n}$ to obtain the uniform superposition
\begin{equation}
  \ket{s} = H^{\otimes n}\ket{0}^{\otimes n} = \frac{1}{\sqrt{N}}\sum_{x}\ket{x},
  \qquad N = 2^n.
\end{equation}

Suppose one basis state $\ket{w}$ is ``marked''. The oracle $O$ acts as
\begin{equation}
  O\ket{x} =
  \begin{cases}
    -\ket{x} & \text{if } x = w, \\
    \phantom{-}\ket{x} & \text{otherwise.}
  \end{cases}
\end{equation}
That is, the oracle flips the phase of the marked state but leaves all other states unchanged.

\subsection*{Diffusion operator as inversion about the average}

Grover's diffusion operator $D$ is defined as
\begin{equation}
  D = 2\ket{s}\bra{s} - I.
\end{equation}

To see its action on a general state
\[
  \ket{\psi} = \sum_{x} a_x\ket{x},
\]
observe that
\[
  \braket{s}{\psi} = \frac{1}{\sqrt{N}}\sum_{x} a_x = \sqrt{N}\,\bar{a},
\]
where $\bar{a}$ is the average amplitude. Then
\begin{align}
  D\ket{\psi} &= \bigl(2\ket{s}\bra{s} - I\bigr)\ket{\psi} \\
  &= 2\ket{s}\braket{s}{\psi} - \ket{\psi} \\
  &= 2\ket{s}\bigl(\sqrt{N}\,\bar{a}\bigr) - \sum_{x} a_x\ket{x} \\
  &= 2\bar{a}\sum_{x}\ket{x} - \sum_{x} a_x\ket{x} \\
  &= \sum_{x}(2\bar{a} - a_x)\ket{x}.
\end{align}
Thus each amplitude $a_x$ is mapped to $2\bar{a} - a_x$: reflection about the mean.

\subsection*{One Grover iteration as interference}

A single Grover iteration applies $G = DO$ to the current state.

\begin{itemize}
  \item The oracle $O$ flips the sign of the marked amplitude $a_w$, leaving all others
    unchanged. This is a phase flip, analogous to the $Z$ gate in the one-qubit example.
  \item The diffusion operator $D$ reflects all amplitudes about their average, which now
    differs from before because of the sign flip.
\end{itemize}

The combined effect is:
\begin{itemize}
  \item The marked state's amplitude is \emph{increased} in magnitude (constructive
    interference).
  \item The unmarked states' amplitudes are slightly \emph{decreased} (destructive interference
    relative to the mean).
\end{itemize}

Repeating $G$ approximately $\frac{\pi}{4}\sqrt{N}$ times rotates the state vector in the
two-dimensional subspace spanned by $\ket{w}$ and the uniform superposition of unmarked states,
bringing it close to $\ket{w}$. Measurement then yields the marked element with high probability.
The speedup arises entirely from carefully controlled interference between the amplitudes.

\section*{Why Interference Matters}

Quantum interference is what enables quantum computers to perform certain computations
exponentially or quadratically faster than classical computers. Classical machines must evaluate
each possibility separately and combine probabilities. Quantum computers evaluate many
possibilities in superposition, then use interference to sculpt the probability landscape so that
correct answers are amplified and incorrect answers are suppressed.

Interference requires the quantum system to maintain coherence. External noise and uncontrolled
interactions cause decoherence, which destroys the delicate interference patterns. Building
scalable, fault-tolerant quantum computers therefore requires keeping coherence long enough and
using error correction to protect the interference structure that algorithms rely on.

\end{document}
